\documentclass[11pt,a4paper]{article}
\usepackage[T1]{fontenc}
\usepackage[utf8]{inputenc}
\usepackage[english]{babel}
\usepackage{amsmath,amssymb}
\usepackage{graphicx}
\usepackage[margin=2cm]{geometry}
\setlength{\parindent}{0pt}
\setlength{\emergencystretch}{3em}
\allowdisplaybreaks

\newsavebox{\fitbox}
\newcommand{\fitmath}[1]{\sbox{\fitbox}{$\displaystyle #1$}%
  \ifdim\wd\fitbox>\linewidth\resizebox{\linewidth}{!}{\usebox{\fitbox}}\else\usebox{\fitbox}\fi}

\title{Proofs}
\author{}
\date{}
\begin{document}
\maketitle

\section*{tests/\allowbreak{}output/\allowbreak{}document\_\allowbreak{}statements.tex}
\begin{enumerate}
\item $\displaystyle \sum_{k=1}^{n} k = \frac{n(n+1)}{2}$

\textbf{Proof.} By induction on $n$.
\begin{itemize}
\item \emph{Base case}, $n = 1$. The left-hand side is $\sum_{k=1}^{1} k = 1$ and the right-hand side is $\frac{1 \cdot \left(1 + 1\right)}{2} = 1$: the two values are equal.
\item \emph{Inductive hypothesis}. Assume the statement for some $n \ge 1$: $\sum_{k=1}^{n} k = \frac{n \cdot \left(n + 1\right)}{2}$.
\item \emph{Inductive step}. We prove the statement for $n + 1$: $\sum_{k=1}^{n + 1} k = \frac{\left(n + 1\right) \cdot \left(n + 1 + 1\right)}{2}$.
\begin{itemize}
\item $\sum_{k=1}^{n + 1} k = \sum_{k=1}^{n} k + \left(n + 1\right)$.
\item By the inductive hypothesis, $\sum_{k=1}^{n} k = \frac{1}{2}\,n + \frac{1}{2}\,n^{2}$.
\item With these replacements, both sides of the statement for $n + 1$ become $1 + \frac{3}{2}\,n + \frac{1}{2}\,n^{2}$.
\end{itemize}
\end{itemize}
$\Box$

\item $\displaystyle \{x \in \mathbb{R} \mid \frac{x - 2}{x + 3} \ge 0\}$

\textbf{Sign table.} Each comparison counts as false where one of its sides has no value. Its truth can change only where one of these expressions in $x$ is zero: $-x + 2$, $x + 3$. At each of these points and on every interval between two consecutive ones the signs are constant, so one value of $x$ decides each.
\begin{itemize}
\item $x < -3$ (for instance $x = -5$): the condition holds.
\item $x = -3$: the condition fails.
\item $-3 < x < 2$ (for instance $x = -2$): the condition fails.
\item $x = 2$: the condition holds.
\item $2 < x$ (for instance $x = 4$): the condition holds.
\end{itemize}
Hence the condition holds exactly for $x < -3 \lor x \ge 2$. $\Box$

\begin{center}
\includegraphics[page=1,width=0.8\textwidth]{mathresearcher_test_plain-plots.pdf}
\end{center}
{\small \textbf{Graph.} Each panel shows the two sides of one comparison, or one side and the axis when the other is $0$. Green bands: where that comparison holds. Gray bands and dashed lines: where a side has no value. Bottom line: where the whole condition holds.}

\item $\displaystyle \forall x \in \mathbb{R}\ \left(\left|x - 1\right| < 2 \iff -1 < x < 3\right)$

\textbf{Sign table.} Each comparison counts as false where one of its sides has no value. Its truth can change only where one of these expressions in $x$ is zero: $x - 3$, $x - 1$, $-x - 1$. At each of these points and on every interval between two consecutive ones the signs are constant, so one value of $x$ decides each.
\begin{itemize}
\item $x < -1$ (for instance $x = -3$): the left side is false, the right side is false.
\item $x = -1$: the left side is false, the right side is false.
\item $-1 < x < 1$ (for instance $x = 0$): the left side is true, the right side is true.
\item $x = 1$: the left side is true, the right side is true.
\item $1 < x < 3$ (for instance $x = 2$): the left side is true, the right side is true.
\item $x = 3$: the left side is false, the right side is false.
\item $3 < x$ (for instance $x = 5$): the left side is false, the right side is false.
\end{itemize}
The condition holds on every row, hence for every value of $x$ in the set. $\Box$

\begin{center}
\includegraphics[page=2,width=0.8\textwidth]{mathresearcher_test_plain-plots.pdf}
\end{center}
{\small \textbf{Graph.} Each panel shows the two sides of one comparison, or one side and the axis when the other is $0$. Green bands: where that comparison holds. Gray bands and dashed lines: where a side has no value. Bottom line: where the whole condition holds.}

\item $\displaystyle \forall x \in \mathbb{R}\ \left(\frac{1}{x} > 0 \iff x \ge 0\right)$

\textbf{Sign table.} Each comparison counts as false where one of its sides has no value. Its truth can change only where one of these expressions in $x$ is zero: $x$, $-x$. At each of these points and on every interval between two consecutive ones the signs are constant, so one value of $x$ decides each.
\begin{itemize}
\item $x < 0$ (for instance $x = -2$): the left side is false, the right side is false.
\item $x = 0$: the left side is false, the right side is true.
\end{itemize}
At this value the condition fails. So the statement is false.

\begin{center}
\includegraphics[page=3,width=0.8\textwidth]{mathresearcher_test_plain-plots.pdf}
\end{center}
{\small \textbf{Graph.} Each panel shows the two sides of one comparison, or one side and the axis when the other is $0$. Green bands: where that comparison holds. Gray bands and dashed lines: where a side has no value. Bottom line: where the whole condition holds.}

\item $\displaystyle \forall x \in \mathbb{R}\ (x + = 3)$

\textbf{No proof.} NON SUPPORTATO: espressione incompleta o simbolo inatteso

\end{enumerate}
\end{document}
